\documentclass[a4paper]{article}
\usepackage{pdfpages, pgffor}
\setlength{\fboxrule}{0.3mm}

\begin{document}
\newwrite\myoutput
\immediate\openout\myoutput=\jobname.options
\immediate\write\myoutput{\unexpanded{\documentclass[10pt, handout, aspectratio=169]{beamer}}} %43 for 4x3, 169 for 16x9, 1610 for 16x10
\immediate\closeout\myoutput
\foreach \n in {1,...,8}{
    \IfFileExists{week\n/week\n.tex}{
        \immediate\write18{texify.exe --pdf  --clean week\n/week\n.tex}
        \includepdf[pages=-,nup=2x2,delta=5pt 30pt,landscape=true,frame=true,column=false,noautoscale=false]{week\n.pdf}
    }
}
\end{document}

\usetheme{Frankfurt}
\usecolortheme{crane}

\usepackage{exscale,latexsym,microtype,amsmath,amssymb,amsfonts,graphicx,natbib,times,booktabs,xstring}

\newcommand{\E}{\ensuremath{{\mathbb E}}} % expected value
\newcommand{\R}{\ensuremath{{\mathbb R}}}
\newcommand{\Var}{\ensuremath{{\mathbb V}}} % variance
\newcommand{\frameit}[2]{\begin{frame}\frametitle{#1}\begin{itemize}#2\end{itemize}\end{frame}}
\def\func#1{\mathop{\rm #1}}
\def\er#1{\emph{\color{red}#1}}
\def\newblock{\hskip .11em plus .33em minus .07em}
\def\limfunc#1{\mathop{\rm #1}}%
\setbeamertemplate{blocks}[rounded][shadow=true]
\date{}

\titlegraphic{\includegraphics[height =0.4in]{../HSLU_Logo_EN_Schwarz_rgb}}
\title{Module 9.3: Time Series Analysis with Python\\Fall Term\space\number\year}
\def\theweek{\StrRight{\jobname}{1}}
\author[Week \theweek]{\textbf{Week \theweek}:}
\AtBeginSection{
\frame{
\frametitle{Outline}
\tableofcontents[currentsection]
}}


\institute{{\Large ARMA Models}}
\begin{document}
\frame{\titlepage}
\begin{frame}%
\frametitle{Outline in Weeks}
\begin{enumerate}
\item Introduction; Descriptive Modeling
\item Returns; Autocorrelation; Stationarity
\item ARMA Models
\item Unit Roots; ARIMA Models
\item Volatility Modeling
\item Value at Risk
\item Cointegration
%\item Panel Data
\end{enumerate}
\end{frame}% 

\section{AR Processes}\subsection*{bla}
\frameit{White Noise and Random Walk}{
\item Last week, we encountered two important stochastic processes: the \er{white noise} process
\begin{align*}
Y_t=U_t,\\
\intertext{and the \er{random walk} process}
Y_t=Y_{t-1}+ U_t.
\end{align*}
\item Observe that we can unify the notation by writing
\[
Y_t=\phi_1 Y_{t-1}+ U_t,
\]
where $\phi_1=0$ produces white noise, and $\phi_1=1$ results in the random walk.
}
\begin{frame}{The AR(1) Process}
\begin{itemize}
\item In fact, there is no reason to restrict the parameter $\phi_1$ to just these two values
\item Allowing arbitrary values results in the \er{autoregressive model of order 1}, or \er{AR(1)} model.
\item Usually, an intercept is also addded. The full model is then
\begin{equation*}
Y_t=\alpha +\phi_1 Y_{t-1}+ U_t, \quad  \text{with $U_t$ white noise.}
\end{equation*}
\item We will see soon that the model is stationary if and only if $-1<\phi_1 <1$.
\item The AR(1) process is a member of a very powerful class of models, the \emph{\color{red}%
autoregressive-moving average} (\er{ARMA}) models, with its special cases
\emph{\color{red}autoregressive} (\er{AR}) and \emph{\color{red}moving average} (\er{MA}) models.

\item Goal: find the right model (for forecasting etc.) by matching the correlogram.
\end{itemize}
\end{frame}
\begin{frame}{Mean  of Stationary AR(1)}
\begin{itemize}

\item A first intuition for the \er{stationarity condition} is obtained if we try to
find the (constant) mean and variance of $Y_t$.
 The mean of $Y_t$ is to be solved from
\begin{equation*}
\E[Y_t] = \alpha + \phi_1 \E[Y_{t-1}] + \E[U_t] = \alpha + \phi_1 \E[Y_{t}],
\end{equation*}
which implies
\begin{equation*}
\E[Y_t] = \frac{\alpha}{1-\phi_1},
\end{equation*}
and this requires $\phi_1 \ne 1$.
\end{itemize}
\end{frame}
\begin{frame}{Variance of Stationary AR(1)}
\begin{itemize}
\item Next, because $\{U_t\}$ is white noise, $U_t$  is uncorrelated with $Y_{t-1}$, so
\begin{equation*}
\mathrm{var}(Y_t) = \phi_1^2 \mathrm{var}(Y_{t-1}) + \mathrm{var}(U_t) + 2
\phi_1 \mathrm{cov}(Y_{t-1},U_{t}) = \phi_1^2 \mathrm{var}(Y_{t}) +
\sigma^2,
\end{equation*}
so that, if and only if $|\phi_1 |<1$,
\begin{equation*}
\mathrm{var}(Y_t) = \frac{\sigma^2}{1-\phi_1^2}.
\end{equation*}

\item Note that $\mathrm{var}(Y_t) > \mathrm{var}(Y_{t-1})$ if $|\phi_1 |\ge 1$, i.e., the variance grows without bounds in that case.
\end{itemize}
\end{frame}



\begin{frame}{ACF / PACF of Stationary AR(1)}
\begin{itemize}
\item It can be shown (see optional exercise) that if $|\phi_1|<1$, then
\begin{equation*}
\tau_k = \frac{\gamma_k}{\gamma_0} = \phi_1 ^k,
\end{equation*}
i.e., if the process is stationary, then the ACF decays exponentially (or \er{geometrically}).
\item The PACF satisfies
\[
\tau_{11} = \phi_1 \quad\mbox{and}\quad \tau_{kk} = 0, k>1;
\]
i.e., it drops to zero after the first lag.
\end{itemize}
\end{frame}
\begin{frame}{Summary: Moments of AR(1)}
\begin{itemize}

\item In summary, we have
obtained the following properties for the stationary AR(1) process:
\begin{block}{Properties of AR(1)}
\begin{itemize}
\item $\E(Y_t) =\dfrac \alpha {1-\phi_1 }$;

\item $\mathrm{var}(Y_t) = \dfrac{\sigma ^2}{1-\phi_1 ^2}$;

\item $\tau _k =\phi_1 ^k,\quad k=1,2,\ldots$;

\item $\tau _{kk} =\left\{
\begin{array}{c}
\phi_1 ,\quad k=1, \\
0,\quad k>1.
\end{array}
\right. $
\end{itemize}
\end{block}
\item Hence, if we find a geometrically declining acf, but a pacf which suddenly
cuts off when $k>1$, then we seem to have an AR(1) process.
\item Try playing with this in the sheet ``AR(1)`` in \texttt{simulation.xlsx}.
\end{itemize}
\end{frame}

\begin{frame}{The Random Walk}
The AR(1) process is non-stationary if:
\begin{itemize}
\item $\phi_1 =1$:
\begin{itemize}

 \item now $Y_t=\alpha +Y_{t-1}+ U_t$: a \emph{%
\color{red} random walk with drift}. It satisfies
\begin{eqnarray*}
\E[Y_t]&=&Y_0+\alpha t, \\
\mathrm{var}(Y_t)&=&\sigma ^2t, \\
\mathrm{corr}(Y_t,Y_{t-k})&=&\sqrt{(t-k)/t}.
\end{eqnarray*}
\item Like for a random walk without drift, the correlogram typically stays close to 1, and decreases slowly,
approximately linearly, while the first difference $\Delta Y_t$ is stationary (see exercises).
\end{itemize}
\item $\phi_1 >1$: this is a so-called \emph{\color{red}explosive} process.
The mean and variance increase very fast (exponentially). We usually do not
consider this because it is unrealistic (at least for long periods of time).
\end{itemize}

\end{frame}
\begin{frame}{AR(p) Models}
\begin{itemize}
\item To fit more complicated ACF patterns, the AR(1) model can be extended to the \er{AR($p$)} model if necessary:
\begin{equation*}
Y_t=\alpha +\phi _1Y_{t-1}+\phi _2Y_{t-2}+\ldots +\phi _pY_{t-p}+ U_t.
\end{equation*}
\item A \emph{necessary} (but not sufficient) \er{condition for stationarity} is
\begin{equation*}
\phi _1+\phi _2+\ldots +\phi _p<1 .
\end{equation*}
\item If the model is stationary, then the mean is $\E[Y_t] = \alpha/(1-\sum_{i=1}^p \phi_i)$.
\item The ACF should gradually approach zero, but not necessarily with a
clear pattern.
\item The PACF satisfies
\begin{equation*}
\tau _{kk}=0,\quad k>p.
\end{equation*}

\item Because all autoregressive models are basically regressions (with lagged
variables as regressors), they can simply be estimated by ordinary
least-squares as usual.
\end{itemize}
\end{frame}

\begin{frame}\frametitle{Example}
Is the AR(2) process
\[
Y_t= .5Y_{t-1}+.5Y_{t-2}+U_t
\]
stationary?
\end{frame}
\section{MA and ARMA Processes}\subsection*{bla}
\begin{frame}{The MA(1) Model}
\begin{itemize}
\item The first-order \emph{\color{red}moving average model} --- MA(1) --- is given by
\begin{equation*}
Y_t=\alpha + U_t + \theta_1 U_{t-1},
\end{equation*}
where $\{U_t\}$ is again a white noise process.
\item This process is stationary for
all values of $\theta_1 $. It can be shown that
\begin{eqnarray*}
\E[Y_t] &=& \alpha, \\
\mathrm{var}(Y_t) &=& \sigma ^2 (1 + \theta_1 ^2), \\
\tau_1 = \mathrm{corr}(Y_t,Y_{t-1}) & = & \theta_1/(1 + \theta_1 ^2).
\end{eqnarray*}
\end{itemize}
\end{frame}
\begin{frame}{ACF and PACF of the MA(1) Model}
\begin{itemize}
\item The most important property is that, for all $k>1$,
\begin{equation*}
\tau_k = \mathrm{corr}(Y_t,Y_{t-k})=0,
\end{equation*}
i.e., the ACF has a \emph{\color{red}cut-off point}; shocks that happened longer than 1 period ago have no
effect on $Y_t$.
\item Also, the PACF of an MA(1) process is $\tau_{kk} =
-(-\theta_1)^k\rightarrow\;\;$\emph{\color{red}geometric decay}.
\item Note that the patterns of the ACF and PACF are \er{reversed} compared to the AR(1) process!
\item We will use this to \er{identify} the correct model from the SACF/SPACF.

\end{itemize}
\end{frame}

\begin{frame}\frametitle{MA($q$) process}
\begin{itemize}
\item
Again we may generalize to the MA($q$) model
\begin{equation*}
Y_t=\alpha + U_t + \theta _1 U_{t-1} + \ldots + \theta _q U_{t-q} .
\end{equation*}
\item Now $\tau _k=0$ for $k>q$, and $\tau_{kk}$ still decays exponentially.
\item This will again be used to identify the correct model from the correlogram.
\end{itemize}
\end{frame}
\frameit{Some more properties of MA($q$) Models }{
\item \er{Estimation}: since the past values of $U_{t}$ are unobserved, we cannot estimate the coefficient $\theta_1 $ by OLS.
Python can still estimate the model, based on other methods.
\item \er{Stationarity}: All MA($q$) processes are stationary, regardless of their parameter values.
}

\begin{frame}{ARMA(1, 1) Process}
\begin{itemize}
\item
The most general class of processes is the mixed autoregressive-moving
average (\er{ARMA}) class of models.
\item The simplest one is the ARMA(1,1):
\begin{equation*}
Y_t= \phi_1 Y_{t-1}+ U_t + \theta_1 U_{t-1}.
\end{equation*}
\item The model is stationary if $|\phi_1 |<1$.% and invertible if $|\theta_1 |<1$.
\item The ACF turns out to be
\begin{eqnarray*}
\tau _1 &=&\frac{(1 + \phi_1 \theta_1 )(\phi_1 + \theta_1 )}{1+\theta_1
^2+2\phi_1 \theta_1 }, \\
\tau _k &=&\phi_1 ^{k-1}\tau _1,\quad k>1 \rightarrow\text{\er{geometric decay}}.
\end{eqnarray*}
\item The PACF is also gradually declining (no cut-off point).
\end{itemize}
\end{frame}

\begin{frame}{ARMA(p, q) Process}
\begin{itemize}
\item The general ARMA(p, q) model is
\begin{equation*}
Y_t= \alpha+\phi_1 Y_{t-1}+\ldots+\phi_p Y_{t-p}  U_t + \theta_1 U_{t-1} + \ldots+\theta_q U_{t-q}.
\end{equation*}
\item Like for the ARMA(1, 1), both its ACF and PACF decay exponentially (provided it is stationary), without a clear pattern.
\item \er{Estimation}: Like pure MA models, ARMA models cannot be estimated by OLS because the past values of $U_{t}$ are unobserved. There are other ways of estimating them though.
\item \er{Stationarity}: The stationarity or otherwise of an ARMA process depends only on the AR component, i.e., a necessary condition is that
\begin{equation*}
\phi _1+\phi _2+\ldots +\phi _p<1 .
\end{equation*}
\end{itemize}
\end{frame}
\begin{frame}{A Theorem}
\begin{itemize}
\item The following theorem explains why ARMA processes are so important:
\begin{block}{Theorem}
\emph{\color{red}Any} stationary process can be represented as an ARMA process.
\end{block}
\item So as long as we can find the right model, we can predict \er{any} stationary process\footnote{And even integrated processes, by taking the first difference and modelling that.}!
\item The properties of the ACF and PACF, summarized on the next slide, will help us do this.
\end{itemize}
\end{frame}

\begin{frame}{Summary of Properties}
\begin{itemize}
\item We have found the following properties of the different processes:
\begin{itemize}
\item AR($p$): geometrically decaying ACF, PACF is zero after $p$ lags;

\item MA($q$): ACF is zero after $q$ lags, geometrically decaying PACF;

\item ARMA($p, q$): geometrically decaying ACF and PACF.
\end{itemize}
\item So we can infer the model type and order from the SACF/SPACF for pure AR and MA models.
\item A full ARMA model would be required if both SACF and PACF decline geometrically, but we won't be able to infer the orders then.
\item The usual procedure is to try an  ARMA(1,1) in that case, and test whether the model needs to be extended.
\item The above is called the \er{Box-Jenkins} approach.
\end{itemize}
\end{frame}

\section{Box-Jenkins Approach}\subsection*{bla}
\begin{frame}{Box-Jenkins}
\begin{itemize}
\item In their 1970 textbook, Box and Jenkins proposed an approach to empirical
ARMA modelling that soon became and still is the predominant approach to
univariate time-series analysis and forecasting. Their procedure consists of
three steps:

\begin{enumerate}
\item \textbf{\color{blue}Identification}. This refers to the problem of
selecting an initial ARMA model, i.e., the choice of the orders $p$
and $q$. This is based on inspection of the
graph and the correlogram of $Y_{t}$. Also,
so-called \emph{\color{red}information criteria} may be used for model
selection.

\item \textbf{\color{blue}Estimation}. The unknown autoregressive and moving
average parameters, as well as the variance $\sigma ^{2}$ of the
disturbances, need to be estimated.

\item \textbf{\color{blue}Diagnostic checking}. A correctly specified model
should not display any autocorrelation in the residuals. Therefore, the main
model check is a test for residual autocorrelation. Also other
misspecification tests (heteroskedasticity, normality, structural change)
may be used.
\end{enumerate}

\item If we find some problem with the model in Step 3, then we return to
Step 1 and go through the cycle again, until the tests indicate no further
problem.
\end{itemize}
\end{frame}
\begin{frame}{Identification: Stationarity}
\begin{itemize}
\item Before even choosing $p$ and $q$, it must be ensured that the data are \er{stationary}. An integrated time series displays
very large autocorrelations, which converge to zero only slowly. In
contrast, stationary time series have an autocorrelation function that
decays exponentially, or is not significantly different from zero after a
few lags.

\item Additional information is available from the graph of the time series.
Stationary time series should display \er{mean-reversion}, i.e., they should
fluctuate around a constant mean (or a linear trend, in case of
trend-stationarity). If a series does not display this property, and behaves
more like a random walk, then it may not be stationary.

\item A formal procedure to test this (and how to proceed in case of non-stationarity) will be introduced later.
\end{itemize}
\end{frame}
\begin{frame}{Identification: $p$ and $q$}
\begin{itemize}
\item The second step is choosing $p$ and $q$.
\item Recall that the correlograms of AR, MA and ARMA processes are characterized thusly:

\begin{itemize}
\item AR($p$): geometrically decaying ACF, PACF is zero after $p$ lags;

\item MA($q$): ACF is zero after $q$ lags, geometrically decaying PACF;

\item ARMA($p, q$): geometrically decaying ACF and PACF.
\end{itemize}

\item If neither ACF nor PACF have a clear cut-off point,
start with an ARMA(1,1), estimate it, and test whether the model
needs to be extended.
\item General goal: \emph{\color{red} parsimony}. Find the smallest possible model that describes the data well.
\end{itemize}
\end{frame}

\begin{frame}{Estimation}
\begin{itemize}
\item See exercises.

\end{itemize}
\end{frame}


\begin{frame}{Diagnostic Testing}
\begin{itemize}
\item
After estimating our favorite ARMA model, we obtain the \er{residuals} $\hat{u}%
_{t}$. These residuals should look like
\er{white noise}.
\item If there is significant
autocorrelation left in $\hat{u}_{t}$, we should
\er{extend} the model. E.g.,, if the residuals of an
AR(1) model look like an MA(1) process, then we might try an ARMA(1,1) model
instead.
\item The most often used test for residual autocorrelation in ARMA models is the
Ljung-Box $Q$-statistic, based on the residuals $\hat{u}_{t}$ instead of the
original time series $y_{t}$.
\item It can be shown that if $\hat{u}_{t}$ is a
residual from an ARMA($p,q$) model, then the $Q$-statistic with $m$
correlations has an approximate $\chi _{m-p-q}^{2}$ distribution under the
null hypothesis. This means that we cannot get $p$-values for the 
first $(p+q)$ $Q$-statistics (see exercises).

\item The tests cannot reliably be used to
find out in which direction the model should be \er{extended}. This means
that we may have to try different alternative specifications before we find
a satisfactory choice.
\end{itemize}
\end{frame}
\begin{frame}{Model Selection Criteria}
\begin{itemize}

\item If more than one specification
passes the diagnostic tests (e.g., both an AR(2) and an ARMA(1,1)), then the decision is often based on the
 \emph{\color{red}Akaike
information criterion} (AIC) and the \emph{\color{red}Schwarz criterion}
(SC; a.k.a.\ the \er{Bayesian} information criterion, BIC):
\begin{eqnarray*}
AIC &=&1+\log 2\pi +\log \left( \frac{1}{n}\sum_{t=1}^{n}\hat{u}%
_{t}^{2}\right) +\frac{2k}{T}, \\
SC &=&1+\log 2\pi +\log \left( \frac{1}{n}\sum_{t=1}^{n}\hat{u}_{t}^{2}\right)
+\frac{k\log T}{T}.
\end{eqnarray*}%
\item $k=$ number of parameters  ($p+q+2$ including intercept and $\sigma^2_u$).
\item Constant $1+\log 2\pi $ is irrelevant in comparisons and sometimes
deleted.
\item We choose the model with the smallest AIC or SC. Typically AIC
leads to higher $p$ and $q$ than SC.
\item The idea is similar to the adjusted $R^2$: the second term is a penalty for including too many parameters. Trade-off between \emph{\color{red}%
goodness of fit} and \emph{\color{red}parsimony}. Smaller models are often better for forecasting.
%Two simple criteria that make this trade-off explicit are the adjusted $%
%R^{2} $ and the residual variance $s^{2}$:%
%\begin{equation*}
%\bar{R}^{2}=1-\frac{\sum_{t=1}^{T}\hat{u}_{t}^{2}/(T-k)}{%
%\sum_{t=1}^{T}(Y_{t}-\bar{Y})^{2}/(T-1)},\qquad s^{2}=\frac{1}{T-k}%
%\sum_{t=1}^{T}\hat{u}_{t}^{2},
%\end{equation*}%
%where $k=p+q+1$ is the total number of parameters (excluding $\sigma ^{2}$%
%).\newpage
\end{itemize}
\end{frame}
\section{Forecasting}\subsection*{bla}
%\begin{frame}{Forecasting Terminology}
%\begin{itemize}
%\item The main purpose of ARIMA models is forecasting.
%\item Denote the forecast of $y_{t+s}$ based on information available at time $t$ as
%$f_{t,s}$.
%
%%\item The (conditional) forecast error variance is $E\left[
%%(y_{t+s}-f_{t,s})^{2}|\mathcal{F} _{t}\right] =\limfunc{var}(y_{t+s}|\mathcal{F} _{t})$%
%%. Usually the expression for this variance does not depend on $\mathcal{F} _{t}$,
%%and therefore is equal to $E\left[ (y_{t+s}-f_{t,s})^{2}\right] $.\newpage
%
%%\item For ARMA models, the one-step-ahead forecast follows from the AR($\infty $)
%%representation:%
%%\begin{equation*}
%%_{t,1}=E(y_{t+1}|\mathcal{F} _{t})=\sum_{i=1}^{\infty }\pi _{i}y_{t+1-i},
%%\end{equation*}%
%%and the forecast error variance is $\limfunc{var}(y_{t+1}|\mathcal{F}
%%_{t})=\sigma ^{2}$.
%
%%\item Multi-step forecasts may be defined recursively, so that $f_{t,2}$ is
%%based on $f_{t,1}$ and $\mathcal{F} _{t}$, and so on. The multi-step forecast
%%error variance may be derived from the MA($\infty $) representation.
%
%%To see how a particular ARIMA model leads to a particular forecast, consider
%%the ARIMA(0,1,1) model%
%%\begin{equation*}
%%\Delta y_{t}=u_{t}+\theta _{1}u_{t-1}.
%%\end{equation*}%
%%Given that $u_{t}=y_{t}-E(y_{t}|\mathcal{F} _{t-1})=y_{t}-f_{t-1,1}$, it follows
%%that%
%%\begin{eqnarray*}
%%E(\Delta y_{t+1}|\mathcal{F} _{t}) &=&E(u_{t+1}|\mathcal{F}
%%_{t})+\theta _{1}E(u_{t}|\mathcal{F} _{t}) \\
%%&=&0+\theta _{1}(y_{t}-f_{t-1,1}),
%%\end{eqnarray*}%
%%which implies%
%%\begin{equation*}
%%f_{t,1}=y_{t}+E(\Delta y_{t+1}|\mathcal{F} _{t})=(1+\theta _{1})y_{t}-\theta _{1}f_{t-1,1}.
%%\end{equation*}%
%%This forecast corresponds to \emph{\color{red}adaptive expectations}, and
%%also to \emph{\color{red}exponential smoothing}. Defining $S_{t}=f_{t,0}$
%%and $\alpha =(1+\theta _{1}) $, it follows that%
%%\begin{equation*}
%%S_{t}=\alpha y_{t}+(1-\alpha )S_{t-1}.
%%\end{equation*}%
%%Here $S_{t}$ is interpreted as a smoothed version of $y_{t}$. This is used,
%%e.g., in RiskMetrics.
%%
%%
%%\newpage
%
%
%
%%Above we listed the main ingredients in forecasting based on ARIMA
%%models. Next we consider some of these ingredients in more detail for
%%the AR(1) model and the random walk.
%\item We distinguish \emph{\color{red}in-sample} and \emph{%
%\color{red}out-of-sample} forecasts.
%\item In-sample means that the same data are forecasted that were used for estimating the model.
%\item Out-of sample forecasting means that the sample is divided into two subsamples, say $1,\ldots,T_1$ (the \emph{\color{red}estimation period}\,) and $T_1+1,\ldots,T$ (the \emph{\color{red}holdout period}\,). The former is used for selecting a model and
%estimating its parameters. Based on this information at time $T_1$, we forecast $%
%y_{T_1+1},\ldots ,y_{T}$.

%\item Alternatively, one may roll the estimation window forward after the first forecast. This way, one always uses the most recent information set $\mathcal{F}_t$ for forecasting $y_t+1$. This is the most relevant for taking a model for a test drive.


%\end{itemize}\end{frame}
%
%Consider again the AR(1) model with constant term,%
%\begin{equation*}
%y_{t}=\mu +\phi y_{t-1}+u_{t},
%\end{equation*}%
%where we assume that the stationarity condition $\left\vert \phi \right\vert
%<1$ is satisfied.
%
%The one-step forecast of $y_{t+1}=\mu +\phi y_{t}+u_{t+1}$ is%
%\begin{eqnarray*}
%f_{t,1} &=&E(y_{t+1}|\mathcal{F} _{t}) \\
%&=&\mu +\phi E(y_{t}|\mathcal{F} _{t})+E(u_{t+1}|\mathcal{F} _{t}) \\
%&=&\mu +\phi y_{t},
%\end{eqnarray*}%
%and the one-step forecast variance is%
%\begin{eqnarray*}
%\sigma _{t,1}^{2} &=&\limfunc{var}\left( y_{t+1}-f_{t,1}|\mathcal{F} _{t}\right)
%\\
%&=&\limfunc{var}(u_{t+1}|\mathcal{F} _{t}) \\
%&=&\sigma ^{2}.
%\end{eqnarray*}%
%\newpage
%
%For $s=2$, we find%
%\begin{eqnarray*}
%f_{t,2} &=&\mu +\phi E(y_{t+1}|\mathcal{F} _{t})+E(u_{t+2}|\mathcal{F} _{t}) \\
%&=&\mu +\phi f_{t,1} \\
%&=&\phi ^{2}y_{t}+\mu (1+\phi ).
%\end{eqnarray*}%
%Therefore, the two-step forecast error may be written as%
%\begin{eqnarray*}
%y_{t+2}-f_{t,2} &=&\mu +\phi y_{t+1}+u_{t+2}-\mu -\phi f_{t,1} \\
%&=&\phi (y_{t+1}-f_{t,1})+u_{t+2}
%\end{eqnarray*}%
%with variance%
%\begin{eqnarray*}
%\sigma _{t,2}^{2} &=&\limfunc{var}(y_{t+2}-f_{t,2}|\mathcal{F} _{t}) \\
%&=&\phi ^{2}\sigma _{t}^{2}+\sigma ^{2} \\
%&=&\sigma ^{2}(1+\phi ^{2}).
%\end{eqnarray*}%
%\newpage
%
%Continuing in the same way, it follows that%
%\begin{eqnarray*}
%f_{t,s} &=&\mu +\phi f_{t,s-1}=\phi ^{s}y_{t}+\mu (1+\phi +\ldots +\phi
%^{s-1}), \\
%\sigma _{t,s}^{2} &=&\sigma ^{2}+\phi ^{2}\sigma _{t,s-1}^{2}=\sigma
%^{2}(1+\phi ^{2}+\ldots +\phi ^{2(s-1)}).
%\end{eqnarray*}%
%As $s$ increases, we see that
%\begin{eqnarray*}
%f_{t,s} &\rightarrow &\frac{\mu }{1-\phi }=E(y_{t}), \\
%\sigma _{t,s}^{2} &\rightarrow &\frac{\sigma ^{2}}{1-\phi ^{2}}=\limfunc{var}%
%(y_{t}).
%\end{eqnarray*}
%
%In other words, the long-run forecast of a stationary AR(1) model is just
%the unconditional mean, and the long-run forecast variance is the
%unconditional variance of $y_{t}$. This is a property of all stationary ARMA
%models, not only the AR(1).\newpage
%
%When $\phi =1$, so that $y_{t}$ is a random walk (with drift)
%and the stationarity
%condition is violated, these properties clearly change. We now have%
%\begin{eqnarray*}
%f_{t,s} &=&\phi ^{s}y_{t}+\mu (1+\phi +\ldots +\phi ^{s-1}) \\
%&=&y_{t}+\mu s.
%\end{eqnarray*}%
%Thus the multi-step forecast is linearly increasing if the drift $\mu $ is
%non-zero, or constant if the drift is zero. For the forecast variance, we
%find%
%\begin{eqnarray*}
%\sigma _{t,s}^{2} &=&\sigma ^{2}(1+\phi ^{2}+\ldots +\phi ^{2(s-1)}) \\
%&=&\sigma ^{2}s,
%\end{eqnarray*}%
%which is therefore also linearly increasing. This means that the width of
%the forecast interval $f_{t,s}\pm 2\sigma _{t,s}$ will keep on increasing as
%$s$ increases. This is a property of all integrated processes, not only the
%random walk.\newpage
\begin{frame}{One-Step Ahead Forecasts}
\begin{itemize}
\item Suppose we want to produce forecasts for the ARMA(1, 1) model
\[
Y_{t+1}=\alpha +\phi _{1}Y_{t}+U_{t+1}+\theta_{1}U_{t}.
\]
\item Strategy: replace parameters with estimates, errors $U_t$ with residuals $\hat{u}_t$, and unobserved future $U_t$ with zero. Hence
\[
\widehat{Y}_{t+1}=\hat\alpha +\hat{\phi }_{1}y_{t}+0+\theta_{1}\hat{u}_{t}.
\]
\item Analogous procedure for general ARMA($p$, $q$) models.
\end{itemize}
\end{frame}
\begin{frame}{Multi-Step Forecasts}
\begin{itemize}
\item When forecasting multiple steps ahead, not only the future $U_t$ are unknown, but also the future $y_t$.
\item Solution: recursive procedure. Forecast $y_{t+1}$ first, then use this to forecast $y_{t+2}$, etc. All future errors are replaced with zero.
%Mathematically,
%\begin{equation*}
%f_{t,s}=\mu +\phi _{1}f_{t,s-1}+\ldots +\phi _{p}f_{t,s-p}+\theta
%_{1}E(u_{t+s-1}|\mathcal{F} _{t})+\ldots +\theta _{q}E(u_{t+s-q}|\mathcal{F} _{t}),
%\end{equation*}%
%where%
%\begin{align*}
%f_{t,s-j}&=\left\{
%\begin{array}{ccc}
%E(y_{t+s-j}|\mathcal{F} _{t}), & \text{if} & j<s, \\
%y_{t+s-j} & \text{if} & j\geq s,%
%\end{array}%
%\right. \intertext{and} E(u_{t+s-j}|\mathcal{F} _{t})&=\left\{
%\begin{array}{ccc}
%0, & \text{if} & j<s, \\
%u_{t+s-j}, & \text{if} & j\geq s.%
%\end{array}%
%\right.
%\end{align*}
%\item For integrated processes, the procedure is to first obtain forecasts
%for $\Delta y_{t+s}$, which is stationary, and then use%
%\begin{equation*}
%E(y_{t+s}|\mathcal{F} _{t})=y_{t}+E(\Delta y_{t+1}|\mathcal{F} _{t})+\ldots +E(\Delta
%y_{t+s}|\mathcal{F} _{t}).
%\end{equation*}%
%Thus the forecasts of $\Delta y_{t+i}$ are accumulated. %Similarly, the
%%forecast errors $\Delta y_{t+i}-\Delta f_{t,i}$ are accumulated, which leads
%%to an expression for the forecast error variance.\newpage
\end{itemize}
\end{frame}

%\begin{frame}{Remarks}
%\begin{itemize}
%
%
%\item EViews distinguishes between \emph{\color{red}static} and
%\emph{\color{red}dynamic} forecasts. This is essentially the same as the
%difference between one-step and multi-step forecasts:
%\textquotedblleft static\textquotedblright\ uses
%information up to time $t-1$ to forecast $y_{t}$, for any $t>T_{1}$; this
%is only possible for in-sample forecasting. \textquotedblleft Dynamic\textquotedblright uses
%only the information up to $T$ (or $T_{1}$) to forecast $T+s$ (or $T_{1}+s$%
%).
%
%\end{itemize}
%\end{frame}
%\begin{frame}[fragile]
%\frametitle{Forecast Evaluation}
%\begin{itemize}
%\item Forecast accuracy is measured by \emph{%
%\color{red}mean square error} (MSE) or \emph{\color{red}mean absolute error}
%(MAE): define the prediction error $e_t\equiv y_{t}-\hat{y}_{t}$. Then\footnote{The notation here differs slightly from Brooks.}
%\begin{equation*}
%MSE=\frac{1}{T-T_{1}}\sum_{t=T_{1}+1}^{T}e_t^{2},\quad MAE=%
%\frac{1}{T-T_{1}}\sum_{t=T_{1}+1}^{T}\left\vert e_t\right\vert
%.
%\end{equation*}%
%\item Forecasts from two competing models with associated prediction errors $e_{1,t}$ and $e_{2,t}$ can be compared by regressing $e_{1,t}^2- e_{2,t}^2$ on a constant. Under the null of equal predictive ability, its $t$ statistic (using HAC standard errors) has an asymptotic standard normal distribution. This is known as \emph{\color{red}Diebold-Mariano test}.
%\item Brooks discusses several other criteria.
%\end{itemize}
%\end{frame}

\section{Epilogue}\subsection*{bla}
\begin{frame}\frametitle{Learning Goals}
Students are able to
\begin{itemize}
\item understand the essential characteristics of AR, MA, and ARMA processes;
\item identify the correct model type from the correlogram;
\item apply diagnostic tests to evaluate a fitted model;
\item choose between competing models using information criteria, and
\item use the final model for forecasting.
\end{itemize}
\end{frame}

\frameit{Homework}{
\item Exercise 3
\item Questions 3, 6--9, 10a, 11a, 11e, 12 from Chapter 6 of Brooks (2019)
}
\end{document} 